%==============================================
% METHODS
%==============================================
\section{Methods}

\subsection{Pipeline Architecture}

GlycoSMILES2BAP operates through four sequential modules (Figure~\ref{fig:pipeline}):

\begin{verbatim}
Input: Glycan notation (IUPAC-condensed / WURCS)
    ↓
[Parser] → Abstract Syntax Tree (AST)
    ↓
[CCD Mapper] → Chemical Component Dictionary codes
    ↓
[BAP Generator] → bondedAtomPairs specification
    ↓
Output: AF3-compatible JSON input
\end{verbatim}

\begin{figure}[htbp]
\centering
% Placeholder for pipeline figure
\fbox{\parbox{0.9\textwidth}{\centering\vspace{2cm}\textbf{Figure 1: Pipeline Architecture}\vspace{2cm}}}
\caption{Pipeline architecture showing the four-stage conversion process from glycan notation to AF3-compatible JSON. The input parsing module accepts IUPAC-condensed, WURCS, or GlycoCT formats and produces an abstract syntax tree. The CCD Mapper converts monosaccharide identifiers to PDB CCD codes with stereochemistry-aware mapping. The BAP Generator produces explicit atom-pair bond specifications.}
\label{fig:pipeline}
\end{figure}

\subsection{Input Parsing Module}

The pipeline accepts three input formats commonly used in glycobiology:

\begin{enumerate}
\item \textbf{IUPAC-condensed}: The standard notation used in publications (e.g., \texttt{Gal(b1-4)GlcNAc})
\item \textbf{WURCS}: Web3 Unique Representation of Carbohydrate Structures
\item \textbf{GlycoCT}: A comprehensive XML-based glycan format
\end{enumerate}

The parser module leverages existing glycan parsing tools (GlyLES, GlycanFormatConverter) to convert input strings into a standardized Abstract Syntax Tree (AST) representation. This AST captures:
\begin{itemize}
\item Monosaccharide identity (Glc, Gal, Man, etc.)
\item Anomeric configuration ($\alpha$/$\beta$)
\item Absolute configuration (D/L)
\item Ring type (pyranose/furanose)
\item Modifications (acetylation, sulfation, etc.)
\item Glycosidic linkage positions
\end{itemize}

\subsection{CCD Mapper Module}

The CCD (Chemical Component Dictionary) provides standardized residue identifiers used by the Protein Data Bank. Our mapper supports 28+ monosaccharide configurations, mapping each (monosaccharide, anomer, absolute configuration) triplet to the appropriate CCD code.

\begin{table}[htbp]
\caption{Key CCD Mappings}
\label{tab:ccd}
\centering
\begin{tabular}{lllll}
\toprule
\textbf{Monosaccharide} & \textbf{Anomer} & \textbf{Config} & \textbf{CCD Code} & \textbf{Anomeric C} \\
\midrule
GlcNAc & $\beta$ & D & NAG & C1 \\
GlcNAc & $\alpha$ & D & A2G & C1 \\
Man & $\alpha$ & D & MAN & C1 \\
Man & $\beta$ & D & BMA & C1 \\
Gal & $\beta$ & D & GAL & C1 \\
Gal & $\alpha$ & D & GLA & C1 \\
Fuc & $\alpha$ & L & FUC & C1 \\
Neu5Ac & $\alpha$ & D & SIA & C2 \\
Neu5Gc & $\alpha$ & D & NGC & C2 \\
Xyl & $\beta$ & D & XYS & C1 \\
\bottomrule
\end{tabular}
\end{table}

\textbf{Key design decisions:}

\begin{enumerate}
\item \textbf{Case-insensitive matching}: All monosaccharide names are normalized to lowercase for robust lookup
\item \textbf{Anomeric position tracking}: Sialic acids (Neu5Ac, Neu5Gc, KDN) use C2 as anomeric carbon (ketoses), while aldoses use C1
\item \textbf{Ring oxygen positions}: Pentoses (Xyl, Ara, Rib) use O4; hexoses use O5; sialic acids use O6
\end{enumerate}

\subsection{BAP Generator Module}

The bondedAtomPairs (BAP) generator converts glycan topology into AF3-compatible bond